\documentclass[11pt]{article}
\usepackage[margin=2.5cm]{geometry}
\usepackage{booktabs}
\usepackage{tabularx}
\usepackage[hidelinks]{hyperref}
\usepackage{times}

\title{Experiments 09 and 10:\\ Two Ideas That Did Not Help}
\author{}
\date{July 31, 2026}

\begin{document}
\maketitle

\section*{What we tested}

After experiment 08, two parts of the pipeline were still hand-made instead of learned:

\begin{enumerate}
  \item \textbf{The readout} --- how the model turns the predictions of the patients in a color
        into one prediction for the whole color. Until now this was a plain average.
  \item \textbf{The graph} --- which patients get connected to which. Until now every patient was
        connected to its 20 most similar patients.
\end{enumerate}

Experiment 09 tried five better readouts. Experiment 10 tried better ways to build the graph.
Everything else was kept exactly as in experiment 08, so the numbers can be compared directly.
Dataset: \texttt{pregnancy\_MiniLM\_L6} (61{,}046 visits, 6{,}590 patients).

\section*{Experiment 09: the readout}

\begin{table}[h]
\centering
\begin{tabularx}{\textwidth}{lXc}
\toprule
Readout & What it does & Test AUC (color) \\
\midrule
\texttt{mean} & Simple average (the old method) & $\mathbf{.798 \pm .004}$ \\
\texttt{gated\_attn} & Learns how much each patient counts & $.796 \pm .005$ \\
\texttt{node\_score\_mean} & Averages the patient scores; \emph{no learning at all} & $.796 \pm .002$ \\
\texttt{pma} & Attention-based pooling & $.795 \pm .004$ \\
\texttt{supernode} & An extra node per color that collects the information & $.795 \pm .009$ \\
\texttt{multistat} & Average, maximum and sum together & $.792 \pm .007$ \\
\bottomrule
\end{tabularx}
\caption{Best setting for each readout, 5 runs each. About 246 settings were tried in total.}
\end{table}

\paragraph{Nothing beat the simple average.} All six numbers sit inside the normal run-to-run noise.
Each new readout got the same amount of tuning as the old one and ended up in the same place.

\paragraph{The most telling row is \texttt{node\_score\_mean}.} It has no learning for the color task
at all --- it just averages the per-patient scores. It matches everything else. So the information
the color task needs is already in the per-patient predictions, and no clever way of combining them
adds anything.

\paragraph{One method was clearly worse.} \texttt{multistat} adds the \emph{sum} of the patient
vectors. Colors contain between 1 and 221 patients, so that sum mostly measures how big the color is,
not what is in it. This confused the model.

\paragraph{Verdict.} Keep the simple average. This part of the pipeline is closed.

\section*{Experiment 10: the graph}

Instead of always connecting each patient to exactly 20 others, we let the number of connections
change: keep a few nearest neighbours for everyone, and add any other patient that is similar enough.
We also tried three different ways of measuring similarity.

\begin{table}[h]
\centering
\begin{tabularx}{\textwidth}{llXc}
\toprule
Model & Graph & Best setting found & Test AUC (color) \\
\midrule
rgat & flexible & avg.\ 62.8 connections & $\mathbf{.800 \pm .003}$ \\
rgat & 20 neighbours & the old graph (avg.\ 31.9) & $.796 \pm .002$ \\
rgcn & flexible & avg.\ 62.1 connections & $.785 \pm .002$ \\
rgcn & 20 neighbours & the old graph (avg.\ 31.9) & $.771 \pm .002$ \\
\bottomrule
\end{tabularx}
\caption{Best setting per graph, 5 runs each.}
\end{table}

\paragraph{The winner is not really a flexible graph.} The best setting keeps almost every candidate
connection: for 77\% of the patients nothing is removed at all. So it is simply a denser graph, not an
adaptive one.

\paragraph{More connections does not mean better.} There is no relationship between the number of
connections and the score. Settings with 9 connections per patient and settings with 63 give the same
result. The way similarity is measured does not matter either.

\paragraph{The small gain is probably not real.} The flexible graph did get slightly better numbers,
but it was also given more than twice as many settings to try ($\sim$95 vs.\ 40). The extra tuning
most likely found a slightly better \emph{model}, not a better graph.

\paragraph{The best model has not improved.} The best result here is $.800 \pm .003$, against
$.798 \pm .006$ in experiment 08. After about 460 settings, we are where we started.

\section*{One measurement worth remembering}

While building experiment 10 we found something more interesting than the experiment itself:
\textbf{all the patient embeddings look almost identical to each other.} Two random patients have a
similarity of $.979$ out of $1$. So the ``20 nearest neighbours'' are chosen by differences in the
fourth decimal place.

We tried two standard fixes (centering and whitening). They do spread the numbers out, and they change
the neighbours almost completely --- but they do not help. In every version, a patient's neighbours
share its label only $1.4$ points above pure chance ($.832$ vs.\ $.818$).

In other words: the neighbour connections barely carry any information about the label. This fits what
experiment 08 already showed --- the entire $+20$ point improvement came from the \emph{same-color}
connections, not from the neighbour connections.

\section*{What to do next}

\begin{enumerate}
  \item \textbf{Quick check first.} Rerun the old 20-neighbour graph with the same number of tuning
        trials as the flexible graph. If the gap disappears, experiment 10 is a clean negative result
        like experiment 09.
  \item \textbf{Then look elsewhere.} The readout is not the problem and the graph shape is not the
        problem. What is left is (a) the patient embeddings themselves, which are very poorly spread
        out, and (b) the same-color connections --- the only part of the pipeline that has ever been
        shown to matter.
\end{enumerate}

\end{document}
